% ID: SEL-D0028-OB-E04
% ESTADO: cotejada. Segunda pasada visual completa y borrador compilado revisados.
% Fuente: Selectividad/matematicas_aplicadas/suplente junio/Examen-Matemáticas Aplicadas a las Ciencias Sociales.pdf
% SHA-256: a10146fce64a17e21450973c88d66d43571619781f76bc345b5a4b3dee67a16b
% Página: 2; opción B; ejercicio original 4.

Se quiere estimar la proporción de estudiantes que asiste de forma regular al cine. Para ello, se toma una muestra aleatoria simple de tamaño 300 y se obtiene que de ellos, 210 acuden con regularidad al cine.
\begin{enumerate}[label=\alph*)]
\item \textbf{(1.75 puntos)} Calcule un intervalo de confianza al 92\% para estimar la proporción de estudiantes que va al cine regularmente. \textquestiondown{}Qué error máximo se cometería si se diera como estimación de dicha proporción 0.7?
\item \textbf{(0.75 puntos)} Con el mismo nivel de confianza, siendo la proporción muestral la misma, si queremos que el error sea menor que 0.02, \textquestiondown{}cuántos alumnos como mínimo hay que elegir en la muestra?
\end{enumerate}
